\setcounter{section}{2}

\Needspace{5\baselineskip}
\hypertarget{ppo-and-gae}{%
\section{PPO and GAE}\label{ppo-and-gae}}

PPO derives from TRPO. TRPO introduced importance sampling and restricted update regions, but solving its constrained optimization problem is complex. It is ``precise but expensive'' and impractical for large models. Constraints can instead be incorporated directly into the objective. Two ways to do this produce two PPO variants.

\Needspace{5\baselineskip}
\hypertarget{i.-ppo-penalty}{%
\subsection{I. PPO-Penalty}\label{i.-ppo-penalty}}

\Needspace{5\baselineskip}
The objective is:

\[
\begin{aligned}
J^{\mathrm{PEN}}(\theta)
&=
\mathbb{E}_t\left[
\frac{\pi_\theta(a_t\mid s_t)}
{\pi_{\theta_{\mathrm{old}}}(a_t\mid s_t)}
A_t
-
\beta\cdot
D_{\mathrm{KL}}\left(
\pi_{\theta_{\mathrm{old}}}(\cdot\mid s_t)
\Vert
\pi_\theta(\cdot\mid s_t)
\right)
\right]
\end{aligned}
\]

\begin{itemize}
\tightlist
\item
  First term: TRPO's surrogate objective, Ratio \(\times\) Advantage, improving the policy.
\item
  Second term: a KL penalty, with \(D_{\mathrm{KL}}\) measuring the difference between old and new policy distributions.
\item
  \(\beta\): the penalty coefficient. Larger \(\beta\) permits less dramatic policy change.
\end{itemize}

The first term is the same as TRPO's: find a new policy maximizing expected reward under that policy. The forward KL term chiefly ensures that regions with high old-policy probability \(\pi_{\mathrm{old}}\), such as common human language patterns and basic world knowledge learned in pretraining, retain high probability under \(\pi_{\mathrm{new}}\). This prevents discarding pretrained knowledge to gain reward and causing ``mode collapse.'' Meanwhile, increasing probabilities of difficult solutions that had low old-policy probability incurs no excessive penalty. This ensures updates remain within a credible ``trust region.''

\Needspace{5\baselineskip}
PPO-Penalty's essence is dynamically adjusting \(\beta\) from actual KL after each update rather than fixing it. Let:

\[
d = D_{\mathrm{KL}}(\pi_{\mathrm{old}},\pi_{\mathrm{new}})
\]

\Needspace{5\baselineskip}
Set target KL \(d_{\mathrm{targ}}\), such as \(0.01\):

\begin{enumerate}
\def\labelenumi{\arabic{enumi}.}
\tightlist
\item
\Needspace{5\baselineskip}
  If \(d < d_{\mathrm{targ}}/1.5\), the change is too small and conservative, so reduce \(\beta\), for example:
\end{enumerate}

\[
\beta \leftarrow \beta/2
\]

This allows a larger next step.

\begin{enumerate}
\def\labelenumi{\arabic{enumi}.}
\setcounter{enumi}{1}
\tightlist
\item
\Needspace{5\baselineskip}
  If \(d > d_{\mathrm{targ}}\times 1.5\), the change is too large and aggressive, so increase \(\beta\), for example:
\end{enumerate}

\[
\beta \leftarrow 2\beta
\]

This penalizes the next update more heavily.

\Needspace{5\baselineskip}
\hypertarget{ii.-ppo-clip-mainstream}{%
\subsection{II. PPO-Clip (Mainstream)}\label{ii.-ppo-clip-mainstream}}

\[
r_t(\theta)=\frac{\pi_\theta(a_t\mid s_t)}{\pi_{\theta_{\mathrm{old}}}(a_t\mid s_t)}
\]

\[
\begin{aligned}
J^{\mathrm{CLIP}}(\theta)
&=
\mathbb{E}_t\left[
\min\left(
r_t(\theta)A_t,\,
\mathrm{clip}(r_t(\theta),1-\epsilon,1+\epsilon)A_t
\right)
\right]
\end{aligned}
\]

\(A_t\) is the action's relative value (discussed below). We seek a new policy maximizing expected value; importance sampling yields the expression above. The derivative of \(J\) with respect to \(\theta\) is the parameter-network update gradient. When \(A_t\) is greater than \(0\), we want larger \(\pi_\theta\), equivalently larger \(r_t\). \(E_t\) denotes data sampled under the old policy.

In practice, first run several steps under \(\pi_{\mathrm{old}}\), then perform several updates. Throughout each round, keep \(\pi_{\mathrm{old}}\) fixed as the data-collection policy (implemented as a direct arithmetic average over past samples), repeating:

\Needspace{5\baselineskip}
When \(A_t>0\), the action is good:

\begin{itemize}
\tightlist
\item
  Originally, larger \(r_t\) is preferred, increasing the action probability.
\item
  Clipping: once \(r_t>1+\epsilon\), this term becomes constant \((1+\epsilon)A_t\), with derivative \(0\).
\item
  Result: policy updating stops, preventing an accidentally high score from raising the action probability without limit.
\end{itemize}

\Needspace{5\baselineskip}
When \(A_t<0\), the action is bad:

\begin{itemize}
\tightlist
\item
  Originally, smaller \(r_t\) is preferred, reducing the action probability.
\item
  Clipping: once \(r_t<1-\epsilon\), the term becomes \((1-\epsilon)A_t\), with derivative \(0\).
\item
  Result: policy updating stops, preventing one negatively rewarded action from excessively disrupting policy structure.
\end{itemize}

Once outside the range, \(J\) is ``flattened,'' and the current update seeking to maximize it ``loses motivation'' and naturally stops.

\Needspace{5\baselineskip}
\hypertarget{iii.-why-does-ppo-clip-perform-better}{%
\subsection{III. Why Does PPO-Clip Perform Better?}\label{iii.-why-does-ppo-clip-perform-better}}

\begin{enumerate}
\def\labelenumi{\arabic{enumi}.}
\item
  Stability: although PPO-Penalty adjusts \(\beta\) dynamically, delayed adjustments or oscillation can destabilize training. PPO-Clip applies hard clipping to individual samples and is extremely robust.
\item
  Computation: PPO-Penalty calculates KL between \(\pi_{\mathrm{new}}\) and \(\pi_{\mathrm{old}}\) (simpler than TRPO's second derivatives but still requiring log-probability differences). Both probabilities change, making computation expensive. PPO-Clip needs only multiplication and comparisons (min, clamp), reducing cost.
\item
  Hyperparameters: fixing PPO-Clip's \(\epsilon\) at \(0.2\) usually suits most environments (Atari, MuJoCo, robots), whereas PPO-Penalty is more sensitive to \(d_{\mathrm{targ}}\) and initial \(\beta\).
\end{enumerate}

\Needspace{5\baselineskip}
\hypertarget{iv.-loss-function}{%
\subsection{IV. Loss Function}\label{iv.-loss-function}}

\[
L_{\mathrm{total}}=-\underbrace{L^{\mathrm{CLIP}}}_{\text{Actor loss}}+c_1\underbrace{L^{\mathrm{VF}}}_{\text{Critic loss}}-c_2\underbrace{S}_{\text{Entropy}}
\]

Note the subtle signs: the optimizer is assumed to perform gradient descent, minimizing the objective.

\begin{itemize}
\tightlist
\item
  \(L^{\mathrm{CLIP}}\): reward should be maximized, so gradient descent requires the negative, \(-L^{\mathrm{CLIP}}\).
\item
  \(L^{\mathrm{VF}}\): prediction error (MSE) should be minimized, so use \(+L^{\mathrm{VF}}\) directly.
\item
  \(S\): entropy should be maximized to encourage exploration, so use \(-S\) for gradient descent.
\end{itemize}

Adding them lets a PyTorch/TensorFlow optimizer optimize all objectives together.

\Needspace{5\baselineskip}
During backpropagation, different loss terms send gradients to corresponding variables in different modules:

\Needspace{10\baselineskip}
\begin{verbatim}
Input image --> [Shared CNN layers] --> Extracted feature vector
                                  |
                                  +-----------------------------+
                                  |                             |
                             [Actor head]                  [Critic head]
                         Action probabilities               Value V
\end{verbatim}

\Needspace{5\baselineskip}
\hypertarget{v.-advantage-a_t-generalized-advantage-estimation-gae}{%
\subsection{\texorpdfstring{V. Advantage \(A_t\): Generalized Advantage Estimation (GAE)}{V. Advantage A\_t: Generalized Advantage Estimation (GAE)}}\label{v.-advantage-a_t-generalized-advantage-estimation-gae}}

\Needspace{5\baselineskip}
Advantage \(A_t\) measures how much better action \(a_t\) is than ``average performance'' in that state. The theoretical definition is:

\[
A_t=Q(s_t,a_t)-V(s_t)
\]

In practice, \(Q\) is unknown and must be estimated from samples.

Case 1: one-step estimation

\Needspace{5\baselineskip}
Looking only one step ahead, action \(a_t\)'s \(Q\) value is estimated by:

\[
r_t+\gamma V(s_{t+1})
\]

\Needspace{5\baselineskip}
The estimated \(A_t\) then equals TD error:

\[
\hat{A}_t^{(1)}=\underbrace{r_t+\gamma V(s_{t+1})}_{\approx Q(s_t,a_t)}-V(s_t)=\delta_t
\]

\Needspace{5\baselineskip}
PPO uses GAE by default to estimate advantages. This ingenious design takes an exponentially weighted moving average of TD errors:

\[
\hat{A}_t^{\mathrm{GAE}}
=
\delta_t+(\gamma\lambda)\delta_{t+1}+(\gamma\lambda)^{2}\delta_{t+2}+\cdots
\]

Here \(\lambda\) is a hyperparameter in \([0,1]\).

\Needspace{5\baselineskip}
\(\lambda\) controls the influence of TD errors:

\Needspace{5\baselineskip}
For \(\lambda=0\), bias is high and variance low:

\[
\hat{A}_t=\delta_t=r_t+\gamma V(s_{t+1})-V(s_t)
\]

\(A_t\) is exactly the current one-step TD error.

\begin{itemize}
\tightlist
\item
  Advantage: low variance, depending only on one random reward \(r_t\).
\item
  Disadvantage: high bias, heavily dependent on critic accuracy \(V(s_{t+1})\). A poorly trained critic gives incorrect estimates.
\end{itemize}

\Needspace{5\baselineskip}
For \(\lambda=1\), there is no bias but high variance:

\[
\hat{A}_t=\sum_{l=0}^{\infty}\gamma^{l}\delta_{t+l}=\left(\sum_{l=0}^{\infty}\gamma^{l}r_{t+l}\right)-V(s_t)
\]

Expanding cancels intermediate \(V\) terms, leaving Monte Carlo returns (actual returns) minus the baseline.

\begin{itemize}
\tightlist
\item
  Advantage: low bias; actual return is the most accurate ``fact.''
\item
  Disadvantage: extremely high variance from accumulating environmental randomness at every step.
\end{itemize}

Which actions produce \(A_t\)? As discussed later, PPO first runs several steps under the old policy, computes advantages for all visited times \(t\), and then updates. Thus \(A_t\) is computed from the actions actually taken during those steps.

\Needspace{5\baselineskip}
\hypertarget{vi.-practical-workflow-such-as-openai-baselines-or-stable-baselines3}{%
\subsection{VI. Practical Workflow (Such as OpenAI Baselines or Stable Baselines3)}\label{vi.-practical-workflow-such-as-openai-baselines-or-stable-baselines3}}

Updating after each step means serial computation and extremely low GPU utilization. Since importance sampling ensures data reuse, we usually work in a ``large loop'': run several steps at once (such as 2048), collect a batch, train and update the policy on it, then discard it.

Initialization: initialize actor \(\pi_\theta\) and critic \(V_\phi\).

\Needspace{5\baselineskip}
Each iteration has four phases:

\textbf{1. Data collection through interaction}

\begin{itemize}
\tightlist
\item
  Run current policy \(\pi_{\theta_{\mathrm{old}}}\) in the environment for \(T\) steps, such as 2048.
\item
\Needspace{5\baselineskip}
  Collect trajectory data:
\end{itemize}

\[
\{s_t,a_t,r_t,s_{t+1},\log\pi_{\mathrm{old}}(a_t\mid s_t)\}
\]

\begin{itemize}
\tightlist
\item
  Use the critic to compute state values \(V(s_t)\).
\end{itemize}

\textbf{2. Advantage computation}

\begin{itemize}
\tightlist
\item
\Needspace{5\baselineskip}
  Compute TD errors:
\end{itemize}

\[
\delta_t=r_t+\gamma V(s_{t+1})-V(s_t)
\]

\begin{itemize}
\tightlist
\item
  Compute GAE \(\hat{A}_t\), a recursive formula balancing bias and variance.
\end{itemize}

Note: values \(V(s_t)\) are usually computed during collection because PPO actors and critics often share lower network layers (CNN/MLP). Choosing \(a_t\) already requires feeding \(s_t\) through the network, so obtaining the critic's \(V(s_t)\) costs almost nothing extra. Without storing it, training would need another pass of \(s_t\) to obtain \(V(s_t)\) for calculating advantages, duplicating computation and wasting GPU resources.

Note: advantages are computed together after the action steps and before policy-gradient updates. GAE weights TD errors and accumulates them from t through 2048.

\textbf{3. Optimization}

\begin{itemize}
\tightlist
\item
  Important: \(\pi_{\theta_{\mathrm{old}}}\) is fixed as the denominator; new \(\pi_\theta\), the numerator, is optimized.
\item
  Shuffle the \(T\) collected samples into minibatches, such as 64 samples each.
\item
  Repeat the epoch loop, for example 10 times.
\end{itemize}

\Needspace{5\baselineskip}
For each minibatch:

\begin{itemize}
\tightlist
\item
\Needspace{5\baselineskip}
  Compute the new probability ratio:
\end{itemize}

\[
r_t(\theta)=\frac{\pi_\theta(a\mid s)}{\pi_{\mathrm{old}}(a\mid s)}
\]

\begin{itemize}
\tightlist
\item
  Compute clipped loss \(L^{\mathrm{CLIP}}\).
\item
\Needspace{5\baselineskip}
  Compute critic value loss:
\end{itemize}

\[
L^{\mathrm{VF}}=\left(V_\phi(s_t)-V_{\mathrm{target}}\right)^{2}
\]

\begin{itemize}
\tightlist
\item
  Compute entropy regularizer \(S\) to encourage exploration.
\item
\Needspace{5\baselineskip}
  Compute total loss:
\end{itemize}

\[
L=-L^{\mathrm{CLIP}}+c_1L^{\mathrm{VF}}-c_2S
\]

\begin{itemize}
\tightlist
\item
  Backpropagate and update \(\theta\) and \(\phi\).
\end{itemize}

\textbf{4. Synchronize the policy}

\begin{itemize}
\tightlist
\item
\Needspace{5\baselineskip}
  After this round, set:
\end{itemize}

\[
\pi_{\theta_{\mathrm{old}}}\leftarrow\pi_\theta
\]

\begin{itemize}
\tightlist
\item
  Clear the data buffer and start the next large loop.
\end{itemize}

\Needspace{5\baselineskip}
Suppose parameters are:

\begin{itemize}
\tightlist
\item
  \texttt{n\_steps}, buffer size: \(2048\).
\item
  \texttt{batch\_size}, minibatch size: \(64\).
\item
  \texttt{n\_epochs}, reuse count: \(10\).
\end{itemize}

\Needspace{5\baselineskip}
The number of backpropagations per ``large loop'' is:

\Needspace{4\baselineskip}
\begin{enumerate}
\def\labelenumi{\arabic{enumi}.}
\tightlist
\item
\Needspace{5\baselineskip}
  Batching: split 2048 samples into:
\end{enumerate}

\[
2048/64=32
\]

minibatches.

\begin{enumerate}
\def\labelenumi{\arabic{enumi}.}
\setcounter{enumi}{1}
\tightlist
\item
  One pass: sequentially feed these 32 minibatches to the GPU, producing 32 backpropagations.
\item
  Repeated reuse: run 10 epochs.
\item
\Needspace{5\baselineskip}
  Total:
\end{enumerate}

\[
32\times 10=320
\]

Thus each large loop produces 320 backpropagations.

\FloatBarrier
\Needspace{5\baselineskip}
\hypertarget{references-58}{%
\subsection{References}\label{references-58}}

\vspace{0.25em}
\begin{itemize}
\tightlist
\item
  Schulman, J., Wolski, F., Dhariwal, P., Radford, A., \& Klimov, O. (2017). \href{https://arxiv.org/abs/1707.06347}{Proximal Policy Optimization Algorithms}. arXiv:1707.06347.
\item
  Schulman, J., Moritz, P., Levine, S., Jordan, M., \& Abbeel, P. (2016). \href{https://arxiv.org/abs/1506.02438}{High-Dimensional Continuous Control Using Generalized Advantage Estimation}. ICLR 2016.
\end{itemize}

\clearpage
